\chapter{Geometric Duality}
\label{ch:geometric_duality}

\section{Point-Line Duality}
\label{sec:duality_point_line}

\index{point-line duality}\index{dual transform}Geometric duality is a
systematic correspondence that swaps the roles of points and lines (in the
plane) or, more generally, of points and hyperplanes (in $\mathbb{R}^d$),
while \emph{preserving incidence and order}. Its value is algorithmic: many
problems that are hard for points become easy for lines, and vice versa, so a
duality transform lets us solve a problem on whichever side of the
correspondence is more convenient. The fundamental planar transform,
developed in Section~\ref{sec:math_duality} of the mathematical foundations,
maps the point $p = (a,b)$ to the non-vertical line
\[
  p^* : \quad y = ax - b,
\]
and maps a non-vertical line $\ell : y = mx + c$ to the point
\[
  \ell^* = (m, -c).
\]
The choice of the \emph{negative} intercept $-b$ is what makes the transform
its own inverse up to the sign convention, so that ``dualizing twice''
returns the original object.

\begin{definition}[Dual Transform]\label{def:geo_duality_transform}
  Let $p = (a,b)$ be a point and $\ell : y = mx + c$ a non-vertical line of
  the plane. The \textbf{dual line} of $p$ is $p^*: y = ax - b$, and the
  \textbf{dual point} of $\ell$ is $\ell^* = (m, -c)$. The map $p \mapsto
  p^*$ on points and $\ell \mapsto \ell^*$ on lines is the \textbf{dual
  transform}. A property that refers to the original plane is said to hold
  in the \textbf{primal}; the transform maps primal objects to
  \textbf{dual} objects.
\end{definition}

\begin{theorem}[Involution]\label{thm:geo_duality_involution}
  The dual transform is an involution on the set of points together with
  non-vertical lines: for every point $p$ and every non-vertical line
  $\ell$,
  \[
    (p^*)^* = p \qquad\text{and}\qquad (\ell^*)^* = \ell .
  \]
\end{theorem}

\begin{proof}
  For a point $p = (a,b)$, the dual line is $p^* : y = ax - b$, i.e., the
  line of slope $a$ and intercept $-b$. The dual of that line is the point
  $(a, -(-b)) = (a,b) = p$. For a line $\ell : y = mx + c$, the dual point
  is $(m, -c)$; its dual line is $y = m x - (-c) = mx + c = \ell$.
\end{proof}

\begin{example}[Basic Dual Pairs]\label{ex:geo_duality_basic}
  The point $p = (2,3)$ has dual line $p^* : y = 2x - 3$; the line
  $\ell : y = -x + 1$ has dual point $\ell^* = (-1, -1)$. Dually, the line
  through the dual points $(2,3)$ and $(-1,-1)$ in the dual plane is
  $y = 2x - 3$ restricted to the two coordinates --- the point $(-1,-1)$
  lies on $p^*$ because $-1 = 2(-1) - 3$, which is exactly the incidence
  property that Section~\ref{sec:duality_incidence} makes precise.
\end{example}

The transform as stated applies only to non-vertical lines and to points of
finite coordinates. Vertical lines $\ell : x = c$ have no finite slope and
are mapped, in the projective closure, to \emph{points at infinity} of the
dual plane; conversely a point with a prescribed dual line of the form
$x = \text{const}$ corresponds, in the primal, to a direction. The standard
workarounds are to rotate the whole configuration so that no relevant line
is vertical, or to use the homogeneous-coordinate machinery of
Section~\ref{sec:math_homogeneous}, where a vertical line is represented by
a vector whose slope component vanishes and the dual point has an infinite
coordinate. In this chapter we assume, unless stated otherwise, that the
input is in general position: no two input points share an $x$-coordinate
and no dual line is vertical.

\begin{remark}[Conventions and Variants]\label{rem:geo_duality_conventions}
  There are several sign conventions for duality. The convention
  $p = (a,b) \mapsto (y = ax - b)$ used here has the property that the dual
  of a \emph{line} stores its slope in the first coordinate and the negative
  of its intercept in the second; interchanging the sign of the intercept
  (e.g., $y = ax + b$) reverses the above/below relations and flips upper
  and lower hulls into one another. A different and equally useful family of
  transforms is the \emph{normal-form} or \emph{polar} duality
  $p = (a,b) \mapsto \{ x : a x + b y = 1 \}$, which treats all lines
  symmetrically at the price of the point $(0,0)$ being sent to infinity;
  the polar duality is the subject of
  Section~\ref{sec:duality_linear_programming}. The choice of convention
  never changes the complexity of the resulting algorithms, only the
  direction in which inequalities and ``above/below'' relations are read.
\end{remark}

\index{dual line}\index{dual point}\index{duality}

\section{Incidence Preservation}
\label{sec:duality_incidence}

The defining property of the dual transform --- the property that makes it a
duality rather than a mere parameterization --- is that it turns
\emph{point-line incidence} into itself.

\begin{definition}[Incidence]\label{def:geo_duality_incidence}
  A point $p$ and a line $\ell$ are \textbf{incident} if $p$ lies on
  $\ell$. Two lines are \textbf{concurrent} if they share a point; a set of
  points is \textbf{collinear} if all of them lie on one line.
\end{definition}

\begin{theorem}[Incidence Preservation]\label{thm:geo_duality_incidence}
  Let $p = (a,b)$ be a point and $\ell : y = mx + c$ a non-vertical line.
  Then
  \[
    p \in \ell \iff \ell^* \in p^* .
  \]
  That is, the primal statement ``the point $p$ lies on the line $\ell$''
  is true if and only if the dual statement ``the point $\ell^*$ lies on the
  line $p^*$'' is true.
\end{theorem}

\begin{proof}
  Both statements are the equation $b = ma + c$. Indeed, $p \in \ell$ means
  $b = ma + c$. The dual point is $\ell^* = (m, -c)$ and the dual line is
  $p^* : y = ax - b$; the incidence $\ell^* \in p^*$ means
  $-c = a m - b$, which rearranges to exactly $b = ma + c$.
\end{proof}

The power of incidence preservation is that it lifts to any configuration
that is described by incidences. Two primal points $p,q$ determine the line
$\overline{pq}$; dually, the two dual lines $p^*, q^*$ intersect at the
single point $\overline{pq}^*$. A primal line through $k$ points becomes a
dual point common to $k$ lines. This is the first of the two great
translation rules of duality:

\begin{corollary}[Collinearity and Concurrence]\label{cor:geo_duality_collinear_concurrent}
  A set $\{p_1, \dots, p_k\}$ of points is collinear in the primal if and
  only if the dual lines $\{p_1^*, \dots, p_k^*\}$ are concurrent in the
  dual. In particular, three points are collinear if and only if their dual
  lines are concurrent, i.e., pass through one common point.
\end{corollary}

\begin{example}[Detecting Collinearity by Hand]\label{ex:geo_duality_collinearity}
  Consider the points $p_1 = (0,0)$, $p_2 = (1,1)$, and $p_3 = (2,2)$.
  Their dual lines are $p_1^* : y = 0$, $p_2^* : y = x - 1$, and
  $p_3^* : y = 2x - 2$. The three dual lines all pass through the point
  $(1,0)$, since $0 = 1 - 1$ and $0 = 2 \cdot 1 - 2$; hence they are
  concurrent, and by the corollary the primal points are collinear. Indeed
  they lie on $y = x$. The dual of that line is $\ell^* = (1, -0) = (1,0)$,
  the concurrence point.
\end{example}

Two further remarks complete the incidence picture. First, the incidence
theorem extends to \emph{segments}: a primal segment $s = pq$ consists of
all points of the segment, so its set of \emph{lines crossing it} is, in
the dual, the union of the dual lines of all points of $s$ --- a
\emph{double wedge} bounded by $p^*$ and $q^*$, with the vertex
$\overline{pq}^*$ at the crossing point of the two boundary lines. Two
primal segments cross if and only if some primal line meets both, which by
duality happens if and only if the two corresponding double wedges overlap.
This translation is the basis of the duality reductions for segment
intersection and for the red--blue problems of
Section~\ref{sec:segint_red_blue}. Second, incidence alone already makes
duality a tool for \emph{verification}: any combinatorial statement that
asserts ``these points lie on these lines'' has an exactly equivalent dual
statement, so a checkable certificate in one plane is automatically a
checkable certificate in the other.

\index{incidence}\index{collinear}\index{concurrent}\index{double wedge}

\section{Order Reversal}
\label{sec:duality_order_reversal}

Incidence is preserved, but \emph{order} is reversed in a precise sense,
and this reversal is the second great translation rule of duality. We need
two order facts: the local order of a point and a line, and the global
order of a chain of points versus a chain of lines.

\begin{lemma}[Point--Line Order]\label{lem:geo_duality_point_line_order}
  Let $p = (a,b)$ be a point and $\ell : y = mx + c$ a non-vertical line.
  The point $p$ lies \emph{above} $\ell$ if and only if the dual point
  $\ell^*$ lies \emph{above} the dual line $p^*$ at $x = m$. Symmetrically,
  $p$ lies below $\ell$ if and only if $\ell^*$ lies below $p^*$.
\end{lemma}

\begin{proof}
  The point $p$ is above $\ell$ when $b > ma + c$. The dual point has
  $y$-coordinate $-c$, and the dual line evaluated at $x = m$ has value
  $am - b$; the point is above the line when $-c > am - b$, which
  rearranges to $b > ma + c$. The ``below'' statement is the same
  inequality reversed.
\end{proof}

The lemma says that, for a fixed pair $(p, \ell)$, the above/below relation
is preserved by duality. The reversal appears once we compare \emph{many}
points against \emph{many} lines, because the sign convention makes the
dual value $a_i m - b_i$ the \emph{negative} of the signed intercept
$b_i - ma_i$.

\begin{theorem}[Order Reversal]\label{thm:geo_duality_order_reversal}
  Let $P$ be a set of points in general position and fix a slope $m$. Sort
  the points of $P$ by the signed distance $b_i - m a_i$ to the line of
  slope $m$ through the origin, smallest first (the point ``lowest'' in the
  direction perpendicular to slope $m$ comes first). Then the dual lines
  $p_i^*$, sorted by their value $a_i m - b_i$ at $x = m$, appear in
  \emph{exactly the reverse order}, largest first. In particular, the point
  of $P$ lowest in the direction of slope $m$ corresponds to the dual line
  attaining the \emph{largest} value at $x = m$, and vice versa.
\end{theorem}

\begin{proof}
  For each point $p_i = (a_i, b_i)$ let $c_i = b_i - m a_i$, the intercept of
  the line of slope $m$ through $p_i$. The value of the dual line at
  $x = m$ is $a_i m - b_i = -c_i$. Sorting the points by $c_i$ increasing
  is therefore the same as sorting the dual values $-c_i$ decreasing, which
  is exactly the claimed reversal.
\end{proof}

The theorem is the mechanism behind every hull--envelope duality. At a fixed
slope $m$, the lowest primal point is dual to the top of the dual
arrangement at $x = m$; letting $m$ range over all slopes, the chain of
lowest points --- the lower convex hull --- traces the upper envelope of the
dual lines, and the upper hull traces the lower envelope. We record the
statement in the form that is used throughout the rest of the book.

\begin{corollary}[Upper Hull and Lower Envelope]\label{cor:geo_duality_hull_envelope}
  For a point set $P$ in general position, a point $p \in P$ is a vertex of
  the \emph{upper} convex hull of $P$ if and only if its dual line $p^*$
  appears on the \emph{lower} envelope of the dual lines $\{q^* : q \in
  P\}$; the upper-hull vertices, read left to right, appear on the lower
  envelope in the reverse order. Symmetrically, the lower-hull vertices,
  read left to right, appear on the \emph{upper} envelope of the dual lines
  in the same order.
\end{corollary}

\begin{proof}
  A point $p$ is a vertex of the upper hull exactly when it is the highest
  point of $P$ in the direction perpendicular to the slope of some
  supporting line of the upper hull. By Theorem~\ref{thm:geo_duality_order_reversal},
  ``highest in that direction'' is dual to ``largest value at $x = m$'',
  which means that $p^*$ is the topmost line at $x = m$, i.e., a segment of
  the lower envelope. The reversal of the order follows because the slopes
  of the supporting lines of the upper hull increase from the leftmost to
  the rightmost vertex, while the envelope is read from left to right. The
  lower-hull statement is symmetric.
\end{proof}

\begin{example}[Envelope Duality on a Triangle]\label{ex:geo_duality_width}
  Let $P = \{(0,0), (1,1), (2,0)\}$. The upper hull is the chain
  $(0,0), (1,1), (2,0)$. The dual lines are $y = 0$, $y = x - 1$, and
  $y = 2x$. The lower envelope of these three lines is $y = 2x$ for
  $x < -1$, $y = x - 1$ for $-1 < x < 1$, and $y = 0$ for $x > 1$, i.e., the
  lines dual to $(2,0), (1,1), (0,0)$ in that order --- exactly the upper
  hull read from right to left, confirming the corollary. The vertical
  distance between the upper and lower envelopes at $x = 0$ is
  $0 - (-1) = 1$, which is the width of the triangle in the vertical
  direction; the general formula is derived in
  Section~\ref{sec:duality_optimization}.
\end{example}

The order-reversal machinery is what connects duality to the theory of
$k$-sets and levels. A \emph{separating line} with exactly $k$ points of $P$
below it defines a $k$-set (Definition~\ref{def:kset}); by the lemma, a line
$\ell$ with $k$ points below it is dual to a point $\ell^*$ that has exactly
$k$ dual lines \emph{below} it, i.e., a point on the \emph{k-level} of the
dual arrangement. The boundary of a $k$-set is therefore the dual of a
$k$-level, and the whole family of $k$-sets of $P$ is encoded in the single
arrangement of the dual lines. Section~\ref{sec:duality_dual_arrangements}
develops this correspondence, and
Section~\ref{sec:arrangement_levels} of the envelopes chapter studies the
levels themselves.

\index{order reversal}\index{upper hull}\index{lower envelope}\index{k-set}

\section{Dual Arrangements}
\label{sec:duality_dual_arrangements}

Given a set $P$ of $n$ points in general position, the dual lines
$\{p^* : p \in P\}$ form an \textbf{arrangement} in the dual plane: the
decomposition of the plane into vertices (intersection points of two dual
lines), edges (the maximal open segments of dual lines between vertices),
and faces (the open cells between consecutive edges). The arrangement is a
single, global object that encodes \emph{all} incidences and order
relations among the points of $P$, and the theory of arrangements developed
in the envelopes chapter (Section~\ref{sec:arrangement_levels} and the
level structure it introduces) applies verbatim.

\begin{definition}[Dual Arrangement]\label{def:geo_duality_dual_arrangement}
  Let $P$ be a set of $n$ points in general position in the primal plane.
  The \textbf{dual arrangement} $A(P) = A(\{p^* : p \in P\})$ is the
  planar arrangement formed by the $n$ dual lines. Each face of $A(P)$
  corresponds to a set of primal lines sharing the same ``profile'': every
  primal line whose dual point lies in a given face has exactly $k$ points
  of $P$ below it, for a value of $k$ that depends only on the face.
\end{definition}

\begin{definition}[Level of the Dual Arrangement]\label{def:geo_duality_level}
  For an integer $k$, the \textbf{k-level} of the dual arrangement is the
  set of points of $A(P)$ that have exactly $k$ dual lines strictly below
  them; it is a polygonal chain made of one segment per dual line. The
  $0$-level is the lower envelope and the $(n-1)$-level is the upper
  envelope of the dual lines (Definition~\ref{def:lower_envelope}).
\end{definition}

The correspondence between the primal configuration and the dual arrangement
can be read off directly:

\begin{itemize}
  \item each dual line $p^*$ is the set of dual points corresponding to all
    primal lines through $p$;
  \item each vertex $p^* \cap q^*$ is the dual point of the line
    $\overline{pq}$ through $p$ and $q$;
  \item the cells above (respectively below) a segment of $p^*$ correspond
    to the primal lines passing above (respectively below) $p$;
  \item the $k$-level separates the cells whose primal lines have $k$
    points below them from the cells whose primal lines have $k+1$ points
    below them;
  \item by Corollary~\ref{cor:geo_duality_collinear_concurrent}, $k$
    collinear points of the primal correspond to a vertex of the dual
    arrangement where $k$ lines meet.
\end{itemize}

The arrangement is the complete \emph{signature} of the point set: two point
sets have combinatorially isomorphic dual arrangements if and only if all
their order and incidence relations agree, which is precisely the notion of
an order type. This makes arrangements the right data structure for
problems in which many points and many lines interact simultaneously.

The complexity and algorithmic questions about arrangements are classical
and have simple answers for lines.

\begin{theorem}[Arrangement Complexity]\label{thm:geo_duality_arrangement_complexity}
  The arrangement of $n$ lines in the plane has at most $n(n-1)/2$
  vertices, $n^2$ edges, and $n(n+1)/2 + 1$ faces; all of these bounds are
  tight for lines in general position. The arrangement can be constructed in
  $O(n^2)$ time and stored in $O(n^2)$ space. The worst-case complexity of
  the $k$-level is $\Theta(n\sqrt{k+1})$.
\end{theorem}

\begin{proof}[Proof sketch]
  Every pair of lines intersects at most once, giving the vertex bound;
  each line is split by the $n-1$ other lines into $n$ edges, giving $n^2$
  edges total; the face count follows from Euler's formula. The $O(n^2)$
  construction is a standard sweep of the arrangement; the zone-based
  refinement of the sweep is the classical algorithm
  (Edelsbrunner--O'Rourke--Seidel \cite{edelsbrunner1986}). The $k$-level
  bound is due to the theory of levels (Theorem~\ref{thm:level_complexity}
  in the envelopes chapter): the upper bound is $\Theta(n\sqrt{k})$ via the
  ``exponential decay'' argument, and point sets for which the level has
  $\Omega(n\sqrt{k})$ vertices exist; the modern treatment appears in
  Edelsbrunner \cite{edelsbrunner1987}.
\end{proof}

\begin{theorem}[Zone of a Line]\label{thm:geo_duality_zone}
  In an arrangement of $n$ lines, the \emph{zone} of any line $\ell$
  --- the set of cells of the arrangement crossed by $\ell$ --- has
  complexity $O(n)$, and the zone can be constructed in $O(n)$ time once
  the arrangement is available.
\end{theorem}

\begin{proof}[Proof sketch]
  Add $\ell$ to the arrangement as the last line. Each edge of the zone is
  bounded by at most two vertices of the original arrangement, and a
  charging argument (each zone vertex is charged to a distinct vertex of the
  arrangement left of it on its level) shows that the zone has at most
  $O(n)$ vertices; hence it has $O(n)$ edges and faces. See
  \cite{edelsbrunner1986}.
\end{proof}

In the dual, the zone theorem has a natural reading: any line of the primal
crossing $k$ of the cells of the dual arrangement corresponds to a line
whose ``order profile'' (the sequence of points above it) changes $O(n)$
times as it moves. The zone theorem is what makes incremental constructions
in arrangements efficient, and it is the structural reason that duality
reductions for problems such as red--blue segment intersection
(Section~\ref{sec:segint_red_blue}) run in near-linear time.

Finally, the dual arrangement viewpoint extends to the paraboloid lifting of
Chapter~\ref{ch:delaunay_triangulations}: the arrangement of the dual lines
of a point set is the 2D analogue of the arrangement of tangent planes to
the paraboloid, and the Voronoi diagram is the projection of the upper
envelope of that tangent-plane arrangement
(Definition~\ref{def:lifting_upper_envelope},
Theorem~\ref{thm:voronoi_upper_hull}, and
Section~\ref{sec:delaunay_lifting}). The dual arrangement is thus the
common ancestor of the envelope, level, and Voronoi structures used
throughout the book.

\index{arrangement}\index{k-level}\index{zone theorem}

\section{Convex Hulls through Duality}
\label{sec:duality_convex_hulls}

\subsection{1. Overview}

The order-reversal theorems of Section~\ref{sec:duality_order_reversal}
turn the convex-hull problem into an envelope problem and back. Because the
convex hull of $n$ points can be computed in $O(n \log n)$ time
(Chapter~\ref{ch:convex_hulls}) and the lower/upper envelope of $n$ lines
can be computed in $O(n \log n)$ time (Chapter~\ref{ch:envelopes}), duality
gives a second family of algorithms for both problems, with the two
complexities coinciding. The same correspondence, lifted one dimension to
the paraboloid, is the mechanism behind the duality of the Delaunay
triangulation and the Voronoi diagram of
Chapter~\ref{ch:delaunay_triangulations}: the ``hull of lifted points'' view
of Section~\ref{sec:delaunay_lifting} is literally a hull--envelope duality
in one dimension higher.

\subsection{2. Definitions}

\begin{definition}[Upper and Lower Hull]\label{def:geo_duality_upper_lower_hull}
  Let $P$ be a point set in general position. The \textbf{upper hull} of
  $P$ is the chain of vertices of $\operatorname{conv}(P)$ visible from
  above, ordered by increasing $x$-coordinate; the \textbf{lower hull} is
  the complementary chain ordered by increasing $x$-coordinate
  (Definition~\ref{def:upper_lower_hull}). Together they form the boundary
  of the convex hull.
\end{definition}

\begin{definition}[Envelope Pair]\label{def:geo_duality_envelope_pair}
  For a family of dual lines $L = \{p^* : p \in P\}$, the \textbf{lower
  envelope} $L_{\text{low}}$ is the pointwise minimum of the lines and the
  \textbf{upper envelope} $L_{\text{up}}$ their pointwise maximum
  (Definitions~\ref{def:lower_envelope} and \ref{def:upper_envelope} of the
  envelopes chapter). Each envelope is a piecewise-linear convex chain.
\end{definition}

\subsection{3. Theory}

The central theorem is the formal version of
Corollary~\ref{cor:geo_duality_hull_envelope}.

\begin{theorem}[Hull--Envelope Duality]\label{thm:geo_duality_upper_hull_lower_envelope}
  Let $P$ be a set of points in general position and let $L = \{p^* : p \in
  P\}$ be their dual lines. Then:
  \begin{itemize}
    \item the vertices of the \emph{upper hull} of $P$, read left to right,
      are exactly the dual lines of the segments of the \emph{lower
      envelope} of $L$, read right to left;
    \item the vertices of the \emph{lower hull} of $P$, read left to right,
      are exactly the dual lines of the segments of the \emph{upper
      envelope} of $L$, read left to right.
  \end{itemize}
  In particular $p$ is an upper-hull vertex if and only if $p^*$ is on the
  lower envelope of the dual lines.
\end{theorem}

\begin{proof}
  A point $p_i$ is an upper-hull vertex exactly when it maximizes the signed
  distance $b_i - m a_i$ for all supporting lines of slope $m$ over an open
  interval of slopes. By Theorem~\ref{thm:geo_duality_order_reversal},
  maximizing $b_i - ma_i$ is equivalent to minimizing the dual value
  $a_i m - b_i$ at $x = m$, so $p_i^*$ is the line attaining the pointwise
  minimum over that interval of slopes --- a segment of the lower envelope.
  The order reversal along the chains follows by the monotonicity of the
  supporting slopes along the upper hull.
\end{proof}

\begin{corollary}[Computational Equivalence]\label{cor:geo_duality_hull_by_duality}
  The convex hull of $n$ points can be computed from the lower and upper
  envelopes of the $n$ dual lines, and the envelopes can be computed from
  the hull of the dual points, with no overhead beyond the $O(n)$ transform.
  Both problems therefore admit algorithms of the same asymptotic
  complexity.
\end{corollary}

\begin{proof}
  Theorem~\ref{thm:geo_duality_upper_hull_lower_envelope} gives both
  hull chains from both envelopes. Conversely, given the hull of the dual
  points $\{p^*\}$ as points, its upper and lower chains dualize back to the
  lower and upper envelopes of the primal dual lines by the same theorem
  applied to the point set $\{p^*\}$.
\end{proof}

The same duality, in one higher dimension, underlies the Delaunay--Voronoi
duality. Lifting each primal point $p = (x,y)$ to
$p^+ = (x, y, x^2 + y^2)$ on the paraboloid, the Delaunay triangulation of
$P$ is the projection of the \emph{lower} hull of the lifted points, and the
Voronoi diagram is dual to it; equivalently the Voronoi diagram is the
projection of the \emph{upper envelope} of the tangent planes to the
paraboloid at the lifted points. This is developed in full in
Section~\ref{sec:delaunay_lifting} of Chapter~\ref{ch:delaunay_triangulations},
where the lifting theorem
(Theorem~\ref{thm:voronoi_upper_hull}) is proved by computing the vertical
gap $\|x - p\|^2$ between the tangent plane and the paraboloid. The hull in
that chapter is a hull in $\mathbb{R}^3$; the duality here is the same
envelope duality as in Theorem~\ref{thm:geo_duality_upper_hull_lower_envelope},
with planes replacing lines.

\subsection{4. Pseudocode}

The algorithm computes the hull of $P$ by computing the two envelopes of the
dual lines and reading off the vertices.

\begin{algorithm}[htbp]
\caption{Convex Hull via Dual Envelopes}
\label{alg:geo_duality_hull_by_envelopes}
\begin{algorithmic}[1]
\Require a set $P$ of $n$ points in general position.
\Ensure the vertices of $\operatorname{conv}(P)$ in counterclockwise order.
\Function{ConvexHullByDuality}{$P$}
  \State $L \gets \emptyset$
  \For{each point $p = (a,b) \in P$}
    \State add the line $y = ax - b$ to $L$ \Comment{dual line $p^*$}
  \EndFor
  \State $\text{low} \gets \text{LowerEnvelope}(L)$ \Comment{recursive divide and conquer}
  \State $\text{high} \gets \text{UpperEnvelope}(L)$ \Comment{negate and repeat}
  \State $H_{\text{up}} \gets $ dual lines of the segments of $\text{low}$,
         read right to left \Comment{upper hull, left to right}
  \State $H_{\text{low}} \gets $ dual lines of the segments of $\text{high}$,
         read left to right \Comment{lower hull, left to right}
  \State \Return $H_{\text{low}}$ followed by $H_{\text{up}}$ minus its
         first and last vertices
\EndFunction
\end{algorithmic}
\end{algorithm}

The subroutines \text{LowerEnvelope} and \text{UpperEnvelope} are the
divide-and-conquer envelope algorithms of Chapter~\ref{ch:envelopes}
(Algorithm~\ref{alg:compute_lower_envelope}); they run in $O(n \log n)$
time. The dual direction --- computing the envelopes from a convex hull ---
uses any hull algorithm of Chapter~\ref{ch:convex_hulls} on the point set
$\{p^*\}$ and reads off the chains as in the proof of
Corollary~\ref{cor:geo_duality_hull_by_duality}.

\subsection{5. Parameter Selection}

\begin{itemize}
\item \textbf{Sign convention}: the transform $y = ax - b$ and its ``upper
  hull to lower envelope'' correspondence are tied together; switching to
  $y = ax + b$ exchanges upper and lower. Pick one convention and keep it
  consistent across the hull and envelope subroutines.
\item \textbf{General position}: points with equal $x$-coordinates produce
  parallel dual lines, whose envelope has vertical segments; either perturb
  the input slightly, or break ties by a consistent rule in the hull and
  envelope algorithms.
\item \textbf{Exact predicates}: the envelope merges and the hull computations
  both rest on orientation and ordering predicates; the robust predicates of
  Chapter~\ref{ch:convex_hulls} and the geometric-predicates chapter
  (Section~\ref{sec:incircle_predicate} and its orientation test) should be
  used.
\end{itemize}

\subsection{6. Complexity}

\begin{itemize}
\item \textbf{Time}: $O(n \log n)$, the cost of the two envelope
  computations (each $O(n \log n)$) or, dually, of the hull computation on
  the dual points. This matches the direct hull algorithms of
  Chapter~\ref{ch:convex_hulls}.
\item \textbf{Space}: $O(n)$ for the envelopes and the resulting hull.
\end{itemize}

\subsection{7. Usages}

\begin{itemize}
\item \textbf{Envelope-based hulls}: when the input is naturally given as
  lines (visibility, kinetic configurations), computing the hull of the dual
  points is the efficient route.
\item \textbf{Delaunay and Voronoi via lifting}: the same duality one
  dimension higher computes the Delaunay triangulation as the lower hull of
  lifted points and the Voronoi diagram as the upper envelope of tangent
  planes (Section~\ref{sec:delaunay_lifting}, Chapter~\ref{ch:delaunay_triangulations});
  the connection to the farthest-point Voronoi diagram is made in
  Section~\ref{sec:farthest_point_voronoi} of Chapter~\ref{ch:voronoi_diagrams}.
\item \textbf{Hull peeling}: onion peeling (Section~\ref{sec:convex_hull_peeling}
  of Chapter~\ref{ch:convex_hulls}) can be run in the dual by peeling levels
  of the arrangement, one level per hull layer.
\end{itemize}

\subsection{8. Testing Methods and Edge Cases}

\begin{itemize}
\item \textbf{Cross-check}: compare the output against a direct hull
  algorithm (e.g., monotone chain) on random and degenerate inputs.
\item \textbf{Collinear and coincident points}: verify the envelope
  routines handle parallel and identical dual lines without producing
  zero-length chain segments.
\item \textbf{Extreme coordinates}: points with large coordinates make the
  dual intercepts large; verify the predicates used in the envelope merge
  are robust.
\end{itemize}

\subsection{9. References}

\begin{itemize}
\item de Berg, M., et al. \cite{deberg2008}, Ch.~8 and Sec.~11.4.
\item Preparata, F.~P., Shamos, M.~I. \cite{preparata1985}, Ch.~3 and
  Ch.~7.
\item Edelsbrunner, H. \cite{edelsbrunner1987}, Ch.~6.
\item Chazelle, B., Guibas, L.~J., Lee, D.~T. \cite{chazelle1985power}.
\end{itemize}

\index{convex hull}\index{envelope}\index{duality}

\section{Half-Plane Intersection}
\label{sec:duality_half_plane_intersection}

\subsection{1. Overview}

The half-plane intersection problem asks for the intersection
$R = \cap_i h_i$ of a set of $n$ half-planes $h_i$; the region $R$ is empty,
or a convex (possibly unbounded) polygon. This is the dual problem to the
convex hull: by the order-reversal theorems of
Section~\ref{sec:duality_order_reversal}, the intersection of the lower
half-planes of the dual lines is a convex region whose boundary is the lower
envelope, and the lower envelope is the dual of the upper hull. Hence every
algorithm for convex hulls yields an algorithm for half-plane intersection
and conversely. The problem is also the geometric form of a linear program:
testing whether $R$ is nonempty is feasibility, and optimizing a linear
function over $R$ is optimization; both are treated in
Section~\ref{sec:duality_linear_programming} and
Chapter~\ref{ch:linear_programming}.

\subsection{2. Definitions}

\begin{definition}[Half-Plane]\label{def:geo_duality_halfplane}
  A \textbf{half-plane} is a set of the form
  \[
    h : \quad y \le mx + c \quad \text{or} \quad y \ge mx + c ,
  \]
  together with its boundary line $\ell : y = mx + c$. The
  \textbf{intersection region} of a family of half-planes is
  $R = \cap_i h_i$. The half-plane $h$ is \emph{lower} if it is of the form
  $y \le mx + c$ and \emph{upper} otherwise; half-planes with vertical
  boundary lines are handled by rotating the input.
\end{definition}

\begin{definition}[Region of Lines above a Point Set]\label{def:geo_duality_above_lines}
  Let $P$ be a set of points and $L = \{p^* : p \in P\}$ the dual lines. The
  region
  \[
    H = \{ (u,v) : v \le a_p u - b_p \text{ for all } (a_p,b_p) \in P \}
  \]
  is the intersection of the lower half-planes of the dual lines. A dual
  point $(u,v)$ belongs to $H$ exactly when the primal line
  $y = ux - v$ has every point of $P$ on or below it.
\end{definition}

\subsection{3. Theory}

The duality between the two problems is exact.

\begin{theorem}[Half-Plane Intersection is Dual to the Convex Hull]\label{thm:geo_duality_hpi}
  Let $P$ be a point set in general position and let $H$ be the intersection
  of the lower half-planes of the dual lines $\{p^*\}$.
  \begin{itemize}
    \item The upper boundary of $H$ is the lower envelope of the dual lines,
      which by Theorem~\ref{thm:geo_duality_upper_hull_lower_envelope} is
      the dual of the upper hull of $P$.
    \item More generally, let $R$ be the intersection of a set of half-planes
      whose boundary lines are in general position. The edges of the lower
      chain of $R$ correspond to the vertices of the \emph{upper} convex hull
      of the dual points of the boundary lines of the lower half-planes; the
      edges of the upper chain of $R$ correspond to the vertices of the
      \emph{lower} convex hull of the dual points of the boundary lines of
      the upper half-planes.
    \item $R$ is bounded if and only if both chains are nonempty and the two
      envelopes cross; $R$ is empty if and only if the two chains do not
      cross.
  \end{itemize}
\end{theorem}

\begin{proof}
  A dual point $(u,v)$ is on the boundary of $H$ when
  $v = \min_i (a_i u - b_i)$, i.e., when $(u,v)$ lies on the lower envelope
  of the dual lines. Theorem~\ref{thm:geo_duality_upper_hull_lower_envelope}
  identifies the lines of the lower envelope with the upper-hull vertices of
  $P$, giving the first statement. For the general case, let $R$ be cut into
  a lower chain (from the lower half-planes) and an upper chain (from the
  upper half-planes). The lower chain is the lower envelope of the boundary
  lines of the lower half-planes; applying the hull--envelope duality to the
  point set $\{\ell_i^*\}$ of the dual points of those boundary lines shows
  the lower envelope is the dual of the upper hull of $\{\ell_i^*\}$, and
  similarly for the upper chain. $R$ is the set of points between the two
  chains over the $x$-interval where the lower chain is below the upper
  chain; the interval is empty exactly when the chains do not cross.
\end{proof}

\begin{corollary}[Hull and Half-Plane Intersection are Equivalent]\label{cor:geo_duality_hpi_hull}
  Computing the convex hull of $n$ points and computing the intersection of
  $n$ half-planes are computationally equivalent: each reduces to the other
  in linear time via the dual transform.
\end{corollary}

\begin{proof}
  By Theorem~\ref{thm:geo_duality_hpi}, the vertices of the intersection of
  half-planes are the duals of the vertices of the convex hull of the dual
  points, and conversely by Theorem~\ref{thm:geo_duality_upper_hull_lower_envelope}.
  Both transforms cost $O(n)$.
\end{proof}

This equivalence is the classical content of Chazelle--Guibas--Lee's study
of ``the power of geometric duality'' \cite{chazelle1985power}: every
algorithm and every lower bound for convex hulls transfers, through duality,
to half-plane intersection, and both problems are LP-feasibility problems in
the sense of Section~\ref{sec:duality_linear_programming}.

\subsection{4. Pseudocode}

\begin{algorithm}[htbp]
\caption{Half-Plane Intersection via the Dual Hull}
\label{alg:geo_duality_halfplane_intersection}
\begin{algorithmic}[1]
\Require $n$ half-planes $h_i$ in general position, each with boundary line
       $\ell_i : y = m_i x + c_i$ and an orientation (lower or upper).
\Ensure the convex region $R = \cap_i h_i$, or a certificate that $R$ is empty.
\Function{IntersectHalfPlanes}{$\{h_i\}$}
  \State $\text{low} \gets \emptyset$;\quad $\text{high} \gets \emptyset$
  \For{each half-plane $h_i$}
    \State $q_i \gets (m_i, -c_i)$ \Comment{dual point of the boundary line}
    \If{$h_i$ is a lower half-plane ($y \le m_i x + c_i$)}
      \State add $q_i$ to $\text{low}$
    \Else
      \State add $q_i$ to $\text{high}$
    \EndIf
  \EndFor
  \State $C_{\text{up}} \gets $ upper convex hull of the points of
         $\text{low}$ \Comment{dual of the lower chain of $R$}
  \State $C_{\text{down}} \gets $ lower convex hull of the points of
         $\text{high}$ \Comment{dual of the upper chain of $R$}
  \State dualize the edges of $C_{\text{up}}$ to the lower chain and the
         edges of $C_{\text{down}}$ to the upper chain of $R$
  \State $R \gets $ the region between the two chains over the $x$-interval
         where lower $\le$ upper
  \If{the interval is empty}
    \State \Return $R = \emptyset$
  \EndIf
  \State clip the chains at their crossing points and \Return $R$
\EndFunction
\end{algorithmic}
\end{algorithm}

The two hull computations are any $O(n \log n)$ hull algorithm of
Chapter~\ref{ch:convex_hulls}. An alternative that avoids duality entirely
is the divide-and-conquer half-plane intersection of
Chapter~\ref{ch:linear_programming} (Section~\ref{sec:halfplane_intersection}),
which merges two convex polygons in linear time and has the same $O(n \log
n)$ bound; the randomized incremental variant is analyzed in
Section~\ref{sec:duality_linear_programming}.

\subsection{5. Parameter Selection}

\begin{itemize}
\item \textbf{Orientation classes}: the algorithm separates lower and upper
  half-planes; half-planes with vertical boundary lines must be rotated out
  of the input first, or handled by a preliminary $y$-sweep.
\item \textbf{Parallel boundary lines}: parallel lines are the degenerate
  case of general position; the hull of the dual points must tolerate
  coincident points (duplicate dual points).
\item \textbf{Boundedness}: to obtain a bounded region, the input must
  contain both lower and upper half-planes in every relevant direction;
  otherwise $R$ is an unbounded slab and the algorithm reports it as such.
\end{itemize}

\subsection{6. Complexity}

\begin{itemize}
\item \textbf{Time}: $O(n \log n)$ worst case, dominated by the two hull
  computations; the merging of the two chains is $O(n)$.
\item \textbf{Space}: $O(n)$.
\item \textbf{Lower bounds}: matching the convex hull, $\Omega(n \log n)$
  comparisons are required to compute the intersection of $n$ half-planes in
  the comparison model (by Corollary~\ref{cor:geo_duality_hpi_hull} and the
  hull lower bounds of Chapter~\ref{ch:convex_hulls}).
\end{itemize}

\subsection{7. Usages}

\begin{itemize}
\item \textbf{Linear programming feasibility}: deciding whether $R \ne
  \emptyset$ and, if so, finding a point of $R$ is the geometric form of LP
  feasibility (Section~\ref{sec:duality_linear_programming} and
  Section~\ref{sec:geometric_linear_programs} of
  Chapter~\ref{ch:linear_programming}).
\item \textbf{Visibility and kernels}: the kernel of a simple polygon is
  the intersection of half-planes defined by its edges; duality transfers
  the hull machinery to kernel computation.
\item \textbf{Smallest enclosing circle}: the smallest enclosing circle of
  a point set can be computed through the farthest-point Voronoi diagram
  (Section~\ref{sec:farthest_point_voronoi}), whose vertices are dual to
  half-plane constraints; the two views meet in the LP formulation of
  Section~\ref{sec:duality_linear_programming}.
\end{itemize}

\subsection{8. Testing Methods and Edge Cases}

\begin{itemize}
\item \textbf{Empty intersection}: construct inputs whose two envelope chains
  do not cross (e.g., half-planes with no common point) and verify the empty
  certificate.
\item \textbf{Unbounded regions}: parallel ``open'' configurations give an
  unbounded slab; verify the algorithm does not emit spurious vertices at
  infinity.
\item \textbf{Cross-validation}: compare the output polygon against the
  divide-and-conquer intersection of Section~\ref{sec:halfplane_intersection}
  on random inputs.
\end{itemize}

\subsection{9. References}

\begin{itemize}
\item de Berg, M., et al. \cite{deberg2008}, Sec.~4.2 and Ch.~8.
\item Preparata, F.~P., Shamos, M.~I. \cite{preparata1985}, Ch.~7.
\item Chazelle, B., Guibas, L.~J., Lee, D.~T. \cite{chazelle1985power}.
\end{itemize}

\index{half-plane intersection}\index{convex region}

\section{Linear Programming through Duality}
\label{sec:duality_linear_programming}

\subsection{1. Overview}

A linear program (LP) asks for a point maximizing (or minimizing) a linear
objective over a set of linear constraints. In the plane, both the
constraints and the objective are geometric: the feasible region is the
intersection of $n$ half-planes, and the optimum of a bounded objective is a
vertex of that region. Duality serves three distinct purposes in this
setting. First, the \emph{primal--dual} correspondence of linear programming
gives certificates: an infeasibility or optimality claim is witnessed by a
convex combination of the constraints, and in the plane these certificates
are exactly the polar-duality statements of the point--line transform.
Second, the randomized incremental algorithms of Seidel solve fixed-dimension
LPs in expected linear time by maintaining the optimal vertex and repairing
it only when a new constraint violates it. Third, the polar (point--hyperplane)
duality extends the planar theory to all dimensions and is the language in
which high-dimensional hulls and their certificates are phrased. This
section develops these three uses; the algebraic LP theory and its dual
program are treated in Chapter~\ref{ch:linear_programming}
(Section~\ref{sec:lp_duality_applications}).

\subsection{2. Definitions}

\begin{definition}[Geometric Linear Program]\label{def:geo_duality_lp}
  Let $H = \{h_1, \dots, h_n\}$ be half-planes and let $c \in
  \mathbb{R}^2$ be an objective direction. The \textbf{2D linear program}
  is
  \[
    \operatorname{minimize}\; c \cdot x \qquad \text{subject to } x \in
    \cap_{i=1}^n h_i .
  \]
  The program is \textbf{infeasible} if the intersection is empty,
  \textbf{unbounded} if the objective is unbounded below over the region, and
  otherwise has an optimal vertex $x^*$ of the feasible polygon.
\end{definition}

\begin{definition}[Polar Body]\label{def:geo_duality_polar}
  Let $K \subseteq \mathbb{R}^d$ be a convex set containing the origin in
  its interior. The \textbf{polar} of $K$ is
  \[
    K^\circ = \{ y \in \mathbb{R}^d : \langle x, y \rangle \le 1 \text{ for
    all } x \in K \} .
  \]
  For a finite set $S = \{s_1, \dots, s_n\}$, the polar of the polytope
  $\operatorname{conv}(S)$ is the intersection of the half-spaces
  $\{ y : \langle s_i, y \rangle \le 1 \}$.
\end{definition}

The polar is the higher-dimensional analogue of the point--line transform:
a point $s$ is mapped to the half-space of the dual hyperplane
$\{y : \langle s, y \rangle = 1\}$ containing the origin. Where the planar
transform swaps points and lines, polarity swaps points and hyperplanes,
and it inverts the face lattice.

\subsection{3. Theory}

The fundamental structural facts are the bipolar theorem, the Farkas
separation statement, and the face-lattice inversion.

\begin{theorem}[Bipolar Theorem]\label{thm:geo_duality_polar_duality}
  Let $K$ be a closed convex set containing the origin. Then
  $(K^\circ)^\circ = K$. If $K$ is a polytope with the origin in its
  interior, then polarity is a bijection between the faces of $K$ and the
  faces of $K^\circ$ that \emph{reverses dimension}: a $k$-dimensional face
  of $K$ corresponds to a $(d-1-k)$-dimensional face of $K^\circ$. In
  particular, the vertices of $K$ correspond to the facets of $K^\circ$.
\end{theorem}

\begin{proof}[Proof sketch]
  The inclusion $K \subseteq (K^\circ)^\circ$ is immediate from the
  definition; the reverse inclusion follows from the hyperplane separation
  theorem (Theorem~\ref{thm:math_separation} of the foundations chapter):
  if $x \notin K$, a strictly separating hyperplane yields a dual vector $y$
  with $\langle x, y \rangle > 1 \ge \langle z, y \rangle$ for all
  $z \in K$, so $y \in K^\circ$ and $x \notin (K^\circ)^\circ$. The
  face-lattice statement follows because a face $F$ of $K$ is the set where
  a supporting hyperplane is tight, and the supporting hyperplane's normal
  is a vertex of $K^\circ$ whose polar facet is the dual of $F$; the
  dimension count is the rank of the tight constraints.
\end{proof}

\begin{theorem}[Zero in a Convex Hull by Polarity]\label{thm:geo_duality_zero_in_conv}
  Let $S = \{s_1, \dots, s_n\} \subseteq \mathbb{R}^d$. Then
  \[
    0 \in \operatorname{conv}(S)
    \iff
    \bigcap_{i=1}^n \{ x : \langle s_i, x \rangle \ge 1 \} = \emptyset .
  \]
  Equivalently, $0 \notin \operatorname{conv}(S)$ if and only if there
  exists a point $x$ with $\langle s_i, x \rangle \ge 1$ for all $i$, i.e.,
  a hyperplane strictly separating $0$ from $\operatorname{conv}(S)$.
\end{theorem}

\begin{proof}
  If $0 \in \operatorname{conv}(S)$ then $0 = \sum_i \lambda_i s_i$ for
  convex weights $\lambda_i \ge 0$, $\sum_i \lambda_i = 1$. For any $x$,
  multiplying $\langle s_i, x \rangle \ge 1$ by $\lambda_i$ and summing
  gives $0 \ge 1$, a contradiction, so the intersection is empty.
  Conversely, if $0 \notin \operatorname{conv}(S)$, the strong separation
  theorem gives a vector $v$ and a scalar $\alpha > 0$ with
  $\langle v, s_i \rangle \ge \alpha$ for all $i$; scaling, the point
  $v/\alpha$ satisfies the system, so the intersection is nonempty.
\end{proof}

The theorem is the certificate machinery behind LP duality: an infeasible
dual constraint system is witnessed by the convex combination (the dual
variables $\lambda_i$), and the dual LP of
Section~\ref{sec:lp_duality_applications} is exactly the search for such a
combination. This is the sense in which polarity is the
higher-dimensional dual transform, and it is the mechanism behind the
high-dimensional hull and polarity algorithms of
Chapter~\ref{ch:convex_polytopes}.

For the algorithmic side, the 2D LP has a simple randomized incremental
solution.

\begin{theorem}[Randomized Incremental LP]\label{thm:geo_duality_seidel}
  The 2D linear program of Definition~\ref{def:geo_duality_lp} can be
  solved in $O(n)$ expected time by Seidel's randomized incremental
  algorithm: insert the half-planes in random order, maintain the optimal
  vertex of the current intersection, and when the new half-plane violates
  the current optimum, solve a one-dimensional subproblem on the new
  boundary line. For any fixed dimension $d$, the algorithm runs in
  $O(d! \, n)$ expected time.
\end{theorem}

\begin{proof}[Proof sketch]
  Maintain the feasible region $R_i = \cap_{j \le i} h_{\pi(j)}$ and its
  optimal vertex $x_i$ for a random permutation $\pi$. If the next
  half-plane $h_{\pi(i)}$ does not contain $x_{i-1}$, the new optimum $x_i$
  must lie on the boundary line of $h_{\pi(i)}$, so the problem reduces to a
  1D LP on that line over the previous constraints --- an interval, whose
  optimum is an endpoint. By backward analysis, the probability that
  $h_{\pi(i)}$ violates $x_{i-1}$ is $O(1/i)$ (for fixed dimension $d$, the
  $d$ constraints binding at the optimum of a random $i$-set are a uniformly
  random $d$-subset, so the new one binds with probability $d/i$). The
  expected total work is therefore $O(\sum_i d/i \cdot i + n) = O(n)$.
\end{proof}

\subsection{4. Pseudocode}

\begin{algorithm}[htbp]
\caption{Seidel's Randomized Incremental 2D LP}
\label{alg:geo_duality_seidel_lp}
\begin{algorithmic}[1]
\Require half-planes $H = \{h_1, \dots, h_n\}$ and an objective direction $c$.
\Ensure an optimal point $x^*$ over $\cap H$, or a certificate of
       infeasibility or unboundedness.
\Function{SeidelLP}{$H$, $c$}
  \State $\pi \gets$ uniformly random permutation of $\{1, \dots, n\}$
  \State $x^* \gets$ an arbitrary feasible point of the initial bounding
         half-planes
  \For{$i = 1, \dots, n$}
    \If{$x^*$ violates $h_{\pi(i)}$}
      \State $x^* \gets \text{OptimumOnLine}(\partial h_{\pi(i)}, c,
             \{h_{\pi(1)}, \dots, h_{\pi(i-1)}\})$
      \Comment{1D LP on the boundary line}
      \If{the 1D subproblem is infeasible}
        \State \Return infeasible
      \EndIf
    \EndIf
  \EndFor
  \State \Return $x^*$
\EndFunction
\Function{OptimumOnLine}{$\ell$, $c$, constraints}
  \State $[u, v] \gets (-\infty, +\infty)$ \Comment{interval of feasible points of $\ell$}
  \For{each half-plane $h$ in constraints}
    \State clip $[u,v]$ by the intersection of $\ell$ with $h$
  \EndFor
  \If{$u > v$}
    \State \Return infeasible
  \EndIf
  \State \Return the point of $[u,v]$ minimizing $c \cdot x$
\EndFunction
\end{algorithmic}
\end{algorithm}

The bounding half-planes of the initialization can be taken as four faraway
constraints enclosing the input; they are also how unboundedness is
detected (the interval $[u,v]$ stays unbounded in the direction of $c$).

\subsection{5. Parameter Selection}

\begin{itemize}
\item \textbf{Feasibility as a subproblem}: when only a point of $R$ is
  needed, take $c$ to be any direction and accept the returned optimum;
  alternatively run the same algorithm with $c = 0$ (any point of the
  intersection is optimal) to test feasibility.
\item \textbf{Degenerate constraints}: a new half-plane containing the
  current optimum is skipped; coincident and parallel boundary lines must
  be deduplicated so that the 1D subproblem sees a clean interval.
\item \textbf{Exact arithmetic}: the predicate ``does $x^*$ violate $h$''
  and the interval clipping use orientation tests; the robust predicates of
  Chapter~\ref{ch:convex_hulls} should be used to avoid near-degenerate
  misclassification.
\end{itemize}

\subsection{6. Complexity}

\begin{itemize}
\item \textbf{Time}: $O(n)$ expected for fixed dimension $d$ (e.g., $d=2$);
  $O(n \log n)$ deterministic by divide-and-conquer half-plane intersection
  (Section~\ref{sec:halfplane_intersection}); the expected number of
  violated constraints is $O(d \log n)$, each costing $O(nd)$ for the 1D
  repair.
\item \textbf{Space}: $O(n)$.
\item \textbf{Certificates}: infeasibility and optimality are certified by
  the convex combinations of Theorem~\ref{thm:geo_duality_zero_in_conv},
  computable from the binding constraints; the certificate size is $O(d)$.
\end{itemize}

\subsection{7. Usages}

\begin{itemize}
\item \textbf{Feasibility and optimization}: the 2D LP solves feasibility
  of half-plane intersection and hence, by
  Theorem~\ref{thm:geo_duality_hpi}, convex-hull inclusion and separation
  queries in the primal (Section~\ref{sec:geometric_linear_programs}).
\item \textbf{Smallest enclosing circle and Chebyshev center}: the smallest
  enclosing circle of a point set is a 2D LP over the (dual) constraints
  of the farthest-point Voronoi diagram
  (Section~\ref{sec:farthest_point_voronoi}), and the largest empty circle
  centered in a region is the corresponding maximization.
\item \textbf{High-dimensional certificates}: polar duality
  (Theorem~\ref{thm:geo_duality_polar_duality}) provides the certificates
  for high-dimensional hulls and for the LP duality treated in
  Chapter~\ref{ch:linear_programming} (Section~\ref{sec:lp_duality_applications}).
\end{itemize}

\subsection{8. Testing Methods and Edge Cases}

\begin{itemize}
\item \textbf{Infeasibility}: verify the algorithm reports infeasible on
  half-planes with empty intersection (e.g., $y \le 0$ and $y \ge 1$).
\item \textbf{Unboundedness}: an objective pointing into an open slab must
  return the unbounded certificate rather than a spurious optimum.
\item \textbf{Degenerate vertices}: when the optimum is a vertex of degree
  $> 2$ (three half-planes through one point), the algorithm must still
  terminate with the correct value.
\end{itemize}

\subsection{9. References}

\begin{itemize}
\item de Berg, M., et al. \cite{deberg2008}, Sec.~4.4 and Ch.~4 (randomized
  LP).
\item Preparata, F.~P., Shamos, M.~I. \cite{preparata1985}, Ch.~4.
\item Matou\v{s}ek, J. \cite{matousek1999} treats the discrepancy and
  approximation sides of the same duality.
\item Chazelle, B., Guibas, L.~J., Lee, D.~T. \cite{chazelle1985power}.
\end{itemize}

\index{linear programming}\index{polar duality}\index{polar body}
\index{separation theorem}

\section{Applications to Optimization}
\label{sec:duality_optimization}

Duality is not merely a reformulation technique: it reduces a family of
optimization problems over point sets to structurally simpler problems over
arrangements and envelopes, for which the near-linear algorithms of
Chapter~\ref{ch:envelopes} apply. This section collects the classical
applications.

\index{collinearity}\index{concurrence}
\paragraph{Detecting collinearities and concurrences.}
By Corollary~\ref{cor:geo_duality_collinear_concurrent}, a set of $k$
collinear primal points is a set of $k$ concurrent dual lines, i.e., a
vertex of the dual arrangement of degree $k$. Counting or detecting
collinear triples among $n$ points therefore reduces to detecting vertices
of the arrangement of the $n$ dual lines at which three lines meet; in the
dual arrangement these are the vertices of the arrangement lying on levels
$k - 1$ to $k - 1$ measured from the bottom (Definition~\ref{def:geo_duality_level}).
The reduction is the subject of the classic ``power of geometric duality''
paper of Chazelle, Guibas, and Lee \cite{chazelle1985power}, which uses
duality to reduce several point-set problems --- including the detection of
collinearities, of intersections, and of certain optimization queries --- to
arrangement and envelope computations. In the algebraic sense, detecting
collinear triples is a 3SUM-hard problem, and the dual formulation exposes
exactly why: a three-line concurrence is the triple meeting condition of an
arrangement vertex, and no subquadratic algorithm is known in the
comparison model.

\index{width}\index{envelope}
\paragraph{Computing envelopes and the width of a point set.}
The order-reversal theorem reduces envelope computation to convex hulls
(Section~\ref{sec:duality_convex_hulls}) and gives duality a direct role in
the width problem.

\begin{theorem}[Width by Envelope Distance]\label{thm:geo_duality_width}
  Let $P$ be a point set, and let $L_{\text{low}}$ and $L_{\text{up}}$ be
  the lower and upper envelopes of the dual lines of $P$. The width of $P$
  in the direction perpendicular to the line of slope $m$ is
  \[
    w(m) = \frac{L_{\text{up}}(m) - L_{\text{low}}(m)}{\sqrt{1 + m^2}} .
  \]
  The minimum of $w(m)$ is attained at an $x$-coordinate of a vertex of the
  envelopes, i.e., at the slope of a line through two hull vertices.
\end{theorem}

\begin{proof}
  The lowest and highest lines of slope $m$ supporting
  $\operatorname{conv}(P)$ have intercepts
  $c_{\min} = \min_i (b_i - m a_i)$ and $c_{\max} = \max_i (b_i - m a_i)$,
  and the width perpendicular to slope $m$ is
  $(c_{\max} - c_{\min})/\sqrt{1 + m^2}$. By
  Theorem~\ref{thm:geo_duality_order_reversal}, the dual values at $x = m$
  are $a_i m - b_i = -c_i$, so $L_{\text{up}}(m) = -c_{\min}$ and
  $L_{\text{low}}(m) = -c_{\max}$, giving the formula. Since both envelopes
  are piecewise linear and the denominator is fixed between consecutive
  arrangement vertices, the function $w$ is minimized on a segment where it
  is the quotient of a linear by a convex function; its minimum over a
  segment is at an endpoint, which is a vertex of the envelopes, dual to a
  line through two hull vertices.
\end{proof}

The theorem is the duality version of the rotating-calipers width
computation: instead of rotating two parallel supporting lines around the
hull, one reads the vertical gap between the two envelopes of the dual
lines and minimizes the scaled gap over the arrangement vertices. This is a
concrete instance of the general phenomenon that envelope algorithms
(Section~\ref{sec:lower_envelopes} and the Davenport--Schinzel theory of
Section~\ref{sec:davenport_schinzel} of Chapter~\ref{ch:envelopes}) solve
optimization problems whose objective is a minimum or maximum of linear
functions.

\index{ham-sandwich}\index{level}
\paragraph{Ham-sandwich cuts and level-set problems.}
The dual of a line that splits a point set is a point on a fixed level of
the dual arrangement. Concretely, a line $\ell$ with exactly $k$ points of
$P$ below it is dual to a point $\ell^*$ with exactly $k$ dual lines below
it, i.e., a point of the $k$-level (Section~\ref{sec:arrangement_levels}
and Definition~\ref{def:geo_duality_level}). A \emph{bisecting} line for a
set of $n$ points (the median in every direction) therefore corresponds to
the \emph{median level} of the dual arrangement, and the ham-sandwich
theorem for two point sets $A$ and $B$ is a statement about the intersection
of two median levels.

\begin{theorem}[Ham-Sandwich via Median Levels]\label{thm:geo_duality_ham_sandwich}
  Let $A$ and $B$ be two finite point sets with $|A|$ and $|B|$ even. The
  median levels of the dual arrangements of $A$ and of $B$ intersect, and a
  point of their intersection is the dual of a line that simultaneously
  bisects $A$ and bisects $B$.
\end{theorem}

\begin{proof}[Proof sketch]
  For each slope $m$, let $\ell_A(m)$ be the line of slope $m$ with
  $|A|/2$ points of $A$ below it; its dual point lies on the median level of
  the $A$-arrangement at $x = m$. The function $f(m) = b_A(m) - b_B(m)$,
  where $b_A(m)$ and $b_B(m)$ are the intercepts of the two bisecting lines,
  changes sign as $m$ runs from $-\infty$ to $+\infty$ (at $m \to \pm
  \infty$ the two bisecting lines are ordered by the projections of $A$ and
  $B$ on the horizontal axis). By continuity, $f$ has a zero, at which the
  bisecting lines coincide; that common line is the ham-sandwich cut, and
  its dual point lies on both median levels.
\end{proof}

\begin{example}[A Ham-Sandwich Cut by Duality]\label{ex:geo_duality_ham_sandwich}
  Let $A = \{(0,0), (2,2)\}$ and $B = \{(0,2), (2,0)\}$. The line $y = 1$
  has one point of $A$ below it ($(0,0)$) and one point of $B$ below it
  ($(2,0)$), so it bisects both sets. Its dual point is $(0,-1)$. The dual
  lines of $A$ are $y = 0$ and $y = 2x - 2$; at $x = 0$ exactly one of them
  ($y = 2x - 2$, value $-2$) lies below $(0,-1)$, so $(0,-1)$ is on the
  median level of $A$. The dual lines of $B$ are $y = -2$ and $y = 2x$;
  at $x = 0$ the line $y = -2$ lies below $(0,-1)$, again exactly one line,
  so $(0,-1)$ is on the median level of $B$ as well --- the dual point is on
  both median levels, exactly as the theorem predicts.
\end{example}

The level-set view also underlies the theory of $k$-sets and discrepancy
studied in Chapter~\ref{ch:envelopes}: the $k$-set boundaries
(Definition~\ref{def:kset}) are dual to the $k$-levels, and the half-plane
discrepancy of a point set (Definition~\ref{def:halfplane_discrepancy}) is
a statement about how far all the levels of the dual arrangement deviate
from being evenly spaced. The discrepancy chapter of Matou\v{s}ek
\cite{matousek1999} develops the deep connections between levels,
discrepancy, and the balance of point sets; the duality here is the
geometric bridge between the two theories.

\index{duality}\index{level-set problem}\index{discrepancy}
\paragraph{The power of geometric duality in optimization.}
The reductions of this section have a common structure: an optimization
problem over a set of points is transformed into a problem over the 
arrangement of the dual lines, where envelope and level algorithms run in
near-linear or near-quadratic time, and where certificates are read off
from local structure (vertices, zone edges, level crossings) rather than
from global searches. Chazelle, Guibas, and Lee \cite{chazelle1985power}
formalized this scheme for a broad family of problems --- collinearity
detection, nearest and farthest neighbor statistics, convex-hull and
intersection queries --- establishing duality as one of the primary
algorithmic paradigms of computational geometry, alongside the
divide-and-conquer and sweeping paradigms of Chapters
\ref{ch:convex_hulls} and \ref{ch:segment_intersection}.

\section{Exercises}
\label{sec:duality_exercises}

\begin{enumerate}
  \item Verify Theorem~\ref{thm:geo_duality_involution} and
    Theorem~\ref{thm:geo_duality_incidence} algebraically for the points
    $(1,2)$, $(3,-1)$ and the line $y = 2x - 1$: compute both duals and
    check all four incidence statements by hand.
  \item Prove Lemma~\ref{lem:geo_duality_point_line_order} in full detail,
    and show that if the sign of the intercept in the transform is changed
    (i.e., $p^* : y = ax + b$), the above/below relation is reversed instead
    of preserved.
  \item Let $P$ be a set of points in general position. Prove that a point
    $p \in P$ is a vertex of the upper hull of $P$ if and only if $p^*$ is on
    the lower envelope of the dual lines, and that the order of the upper-hull
    vertices is the reverse of the order of the envelope segments
    (Corollary~\ref{cor:geo_duality_hull_envelope}).
  \item Use Theorem~\ref{thm:geo_duality_zero_in_conv} to show that the
    point $(0,0)$ lies in the convex hull of $S = \{(1,0),(0,1),(-1,-1)\}$
    but not in the convex hull of $S' = \{(1,0),(0,1)\}$, by writing down the
    dual systems and checking infeasibility directly.
  \item Describe how the width of a point set is computed from the two
    envelopes of the dual lines, and prove that the minimizing slope
    corresponds to a vertex of the envelopes
    (Theorem~\ref{thm:geo_duality_width}). Compute the width of the triangle
    $\{(0,0),(1,1),(2,0)\}$ this way.
  \item Show that two segments in the primal intersect if and only if their
    dual double wedges overlap, and explain how this reduces red--blue
    segment intersection (Section~\ref{sec:segint_red_blue}) to an
    arrangement or envelope problem in the dual plane.
  \item Implement Algorithm~\ref{alg:geo_duality_halfplane_intersection}
    (half-plane intersection via the dual hull) and cross-validate it against
    the divide-and-conquer intersection of
    Section~\ref{sec:halfplane_intersection} on random and degenerate
    inputs, including empty and unbounded cases.
  \item Verify the ham-sandwich statement of
    Theorem~\ref{thm:geo_duality_ham_sandwich} computationally: for random
    even-size sets $A$ and $B$, compute the median levels of the two dual
    arrangements, find an intersection point, dualize it back, and check that
    the resulting line bisects both sets.
\end{enumerate}

\index{exercises}
